A permutation of $\mathbb{N}$ is a bijection $\pi:\mathbb{N}\to\mathbb{N}$ viewed as a sequence $(\pi(1),\pi(2),\dots)$.
A \emph{monotone $k$-term arithmetic progression in $\pi$} means there exist indices $i_1<\dots<i_k$ such that
$\pi(i_1),\dots,\pi(i_k)$ form an arithmetic progression and are monotone in that index order.
Question: Must every permutation of $\mathbb{N}$ contain a monotone 4-term arithmetic progression?
